% Shared metadata and theorem notation for a FormalizedPapers project.
\newtheorem{theorem}{Theorem}[chapter]
\newtheorem{proposition}[theorem]{Proposition}
\newtheorem{lemma}[theorem]{Lemma}
\newtheorem{corollary}[theorem]{Corollary}
\newtheorem{conjecture}[theorem]{Conjecture}
\theoremstyle{definition}
\newtheorem{definition}[theorem]{Definition}
\newtheorem{remark}[theorem]{Remark}
\newtheorem{notation}[theorem]{Notation}
\newtheorem{convention}[theorem]{Convention}

\providecommand{\label}[1]{}
\providecommand{\source}[1]{}
\providecommand{\uses}[1]{}
\providecommand{\lean}[1]{}
\providecommand{\leanok}{}
\providecommand{\mathlibok}{}
\providecommand{\notready}{}

\newcommand{\Spec}{\operatorname{Spec}}
\newcommand{\Pic}{\operatorname{Pic}}
\newcommand{\Jac}{\operatorname{Jac}}
\newcommand{\GL}{\operatorname{GL}}
\newcommand{\Hom}{\operatorname{Hom}}
\newcommand{\End}{\operatorname{End}}
\newcommand{\C}{\mathbb{C}}
\newcommand{\Q}{\mathbb{Q}}
\newcommand{\Z}{\mathbb{Z}}
\newcommand{\F}{\mathbb{F}}
\newcommand{\P}{\mathbb{P}}
\newcommand{\A}{\mathbb{A}}
\newcommand{\M}{\mathcal{M}}
\newcommand{\T}{\mathbb{G}_m}
\newcommand{\K}{\mathcal{K}}



\chapter*{About This Blueprint}
\label{standard-conjectures-abelian-fourfolds-about}

This project follows \emph{Standard conjectures for abelian fourfolds} by Giuseppe Ancona (2020). The
retrieved source is archived under
\href{../../../../../references/standard-conjectures-abelian-fourfolds/}{\texttt{standard-conjectures-abelian-fourfolds}} and is indexed in
the workspace reference manifest. The corresponding preprint is
\href{https://arxiv.org/abs/1806.03216}{arXiv:1806.03216}.

We work with an abelian variety over a field, its algebraic cycles modulo numerical equivalence, and the intersection pairings induced by a polarization. The fourfold case is the central dimension. The paper decomposes the motive of an abelian variety, analyzes Lefschetz and Hodge--Riemann pairings, and compares Betti, ℓ-adic, and $p$-adic realizations through CM structures and quadratic forms.

This blueprint is an independently worded, source-faithful mathematical
outline of the paper's definitions, constructions, and principal results. It
should be read alongside the original source; the source anchor on each node
points to the complete extracted TeX. Statements are intentionally marked
\notready until a matching Lean declaration and proof have been added.

\chapter{Foundations and setup}
\label{chap:standard-conjectures-abelian-fourfolds-foundations}

\section{Standing hypotheses}
\label{sec:standard-conjectures-abelian-fourfolds-standing}

We work with an abelian variety over a field, its algebraic cycles modulo numerical equivalence, and the intersection pairings induced by a polarization. The fourfold case is the central dimension.

\begin{convention}[Standing hypotheses]
\label{standard-conjectures-abelian-fourfolds-setup}
\source{standard-conjectures-abelian-fourfolds}
The objects and hypotheses described above are in force throughout the
formalization of this paper.
\notready
\end{convention}

\chapter{Constructions and proof architecture}
\label{chap:standard-conjectures-abelian-fourfolds-method}

\section{Geometric and cohomological method}
\label{sec:standard-conjectures-abelian-fourfolds-method}

The paper decomposes the motive of an abelian variety, analyzes Lefschetz and Hodge--Riemann pairings, and compares Betti, ℓ-adic, and $p$-adic realizations through CM structures and quadratic forms. The formalization should preserve these intermediate constructions
as separate dependency nodes before addressing the principal result.

\begin{definition}[Central object of the paper]
\label{standard-conjectures-abelian-fourfolds-central-object}
\source{standard-conjectures-abelian-fourfolds}
The central moduli, cycle, representation, height, or filtration object
introduced in the paper is taken with the hypotheses in the standing
convention.
\uses{standard-conjectures-abelian-fourfolds-setup}
\notready
\end{definition}

\chapter{Principal results}
\label{chap:standard-conjectures-abelian-fourfolds-results}

\section{Main theorem or conjecture}
\label{sec:standard-conjectures-abelian-fourfolds-results}

\begin{theorem}[Hodge standard conjecture in dimension four]
\label{standard-conjectures-abelian-fourfolds-main}
\source{standard-conjectures-abelian-fourfolds}
\uses{standard-conjectures-abelian-fourfolds-setup}
The standard conjecture of Hodge type holds for abelian fourfolds over fields of positive characteristic.
\notready
\end{theorem}

\begin{theorem}[Numerical and ℓ-adic homological equivalence]
\label{standard-conjectures-abelian-fourfolds-secondary}
\source{standard-conjectures-abelian-fourfolds}
\uses{standard-conjectures-abelian-fourfolds-main}
For an abelian fourfold, numerical equivalence agrees with ℓ-adic homological equivalence for infinitely many prime numbers \(\ell\).
\notready
\end{theorem}

